\textbf{Phần 1. Trắc nghiệm lựa chọn}
\setcounterpageref{ex}{0}

\begin{ex} %[3X1V73]
Bảng thống kê số cây trồng của 4 tổ: Tổ 1: 18 cây; Tổ 2: 24 cây; Tổ 3: 16 cây; Tổ 4: 22 cây. Tổ trồng được nhiều cây nhất là
\choice
{Tổ 1}
{Tổ 4}
{\True Tổ 2}
{Tổ 3}
\loigiai{
Tổ 2 trồng được 24 cây, nhiều nhất trong 4 tổ.
\textbf{Đáp án của câu hỏi là C}
}
\end{ex}

\begin{ex} %[3X1V73]
Tổng số cây trồng của cả 4 tổ ở câu trên là
\choice
{\True $80$ cây}
{$75$ cây}
{$85$ cây}
{$90$ cây}
\loigiai{
Tổng số cây: $18 + 24 + 16 + 22 = 80$ cây.
\textbf{Đáp án của câu hỏi là A}
}
\end{ex}

\begin{ex} %[3X1V73]
Tổ trồng nhiều cây nhất nhiều hơn tổ trồng ít cây nhất bao nhiêu cây?
\choice
{$6$ cây}
{$7$ cây}
{$9$ cây}
{\True $8$ cây}
\loigiai{
Nhiều nhất là Tổ 2 (24 cây), ít nhất là Tổ 3 (16 cây). Chênh lệch: $24 - 16 = 8$ cây.
\textbf{Đáp án của câu hỏi là D}
}
\end{ex}

\begin{ex} %[3X1V73]
Số học sinh giỏi của ba lớp: Lớp 3A có 15 bạn, Lớp 3B có 18 bạn, Lớp 3C có 12 bạn. Khẳng định nào sau đây là đúng?
\choice
{Lớp 3C có nhiều học sinh giỏi nhất}
{\True Lớp 3B có nhiều học sinh giỏi nhất}
{Lớp 3A có ít học sinh giỏi nhất}
{Cả ba lớp có 40 học sinh giỏi}
\loigiai{
Lớp 3B có 18 bạn học sinh giỏi, nhiều nhất.
\textbf{Đáp án của câu hỏi là B}
}
\end{ex}

\begin{ex} %[3X1V73]
Thống kê số điểm 10 trong tuần của Mai: Thứ Hai: 3; Thứ Ba: 2; Thứ Tư: 4; Thứ Năm: 1; Thứ Sáu: 2. Ngày Mai đạt nhiều điểm 10 nhất là
\choice
{Thứ Hai}
{Thứ Sáu}
{\True Thứ Tư}
{Thứ Ba}
\loigiai{
Thứ Tư Mai đạt 4 điểm 10, nhiều nhất.
\textbf{Đáp án của câu hỏi là C}
}
\end{ex}

\begin{ex} %[3X1V73]
Tổng số điểm 10 Mai đạt được trong cả tuần là
\choice
{\True $12$ điểm $10$}
{$10$ điểm $10$}
{$11$ điểm $10$}
{$14$ điểm $10$}
\loigiai{
Tổng số điểm 10: $3 + 2 + 4 + 1 + 2 = 12$ điểm 10.
\textbf{Đáp án của câu hỏi là A}
}
\end{ex}

\begin{ex} %[3X1V73]
Một cửa hàng bán được trong 3 ngày: Ngày 1: $120\text{ kg}$; Ngày 2 gấp đôi Ngày 1; Ngày 3 ít hơn Ngày 2 là $50\text{ kg}$. Ngày 3 bán được bao nhiêu ki-lô-gam?
\choice
{$240\text{ kg}$}
{$200\text{ kg}$}
{$180\text{ kg}$}
{\True $190\text{ kg}$}
\loigiai{
Ngày 2: $120 \times 2 = 240\text{ kg}$. Ngày 3: $240 - 50 = 190\text{ (kg)}$.
\textbf{Đáp án của câu hỏi là D}
}
\end{ex}

\begin{ex} %[3X1V73]
Số lượng học sinh đăng ký các câu lạc bộ: Bơi lội: 45 bạn; Cờ vua: 30 bạn; Bóng đá: 60 bạn; Cầu lông: 25 bạn. Câu lạc bộ có ít bạn đăng ký nhất là
\choice
{Cờ vua}
{\True Cầu lông}
{Bơi lội}
{Bóng đá}
\loigiai{
Cầu lông có 25 bạn đăng ký, ít nhất.
\textbf{Đáp án của câu hỏi là B}
}
\end{ex}

\begin{ex} %[3X1V73]
Số bạn đăng ký Bóng đá gấp mấy lần số bạn đăng ký Cờ vua?
\choice
{$3$ lần}
{$1$ lần}
{\True $2$ lần}
{$4$ lần}
\loigiai{
$60 : 30 = 2$ lần.
\textbf{Đáp án của câu hỏi là C}
}
\end{ex}

\begin{ex} %[3X1V73]
Bảng thống kê khối lượng thu hoạch lúa của 3 thửa ruộng: Thửa A: $850\text{ kg}$; Thửa B: $920\text{ kg}$; Thửa C: $780\text{ kg}$. Tổng khối lượng lúa cả ba thửa ruộng là
\choice
{\True $2550\text{ kg}$}
{$2500\text{ kg}$}
{$2600\text{ kg}$}
{$2450\text{ kg}$}
\loigiai{
$850 + 920 + 780 = 2550\text{ (kg)}$.
\textbf{Đáp án của câu hỏi là A}
}
\end{ex}

\begin{ex} %[3X1V73]
Để kiểm tra sĩ số và tình hình chuyên cần của học sinh hàng ngày, cán bộ lớp thường sử dụng
\choice
{Thước đo}
{Nhiệt kế}
{Compa}
{\True Sổ điểm danh (bảng theo dõi số liệu)}
\loigiai{
Sổ theo dõi/điểm danh là công cụ thống kê số liệu hàng ngày.
\textbf{Đáp án của câu hỏi là D}
}
\end{ex}

\begin{ex} %[3X1V73]
Trong một bài khảo sát về con vật nuôi yêu thích của 100 học sinh: Chó: 40 bạn; Mèo: 35 bạn; Cá cảnh: 15 bạn; Chim cảnh: 10 bạn. Số bạn thích Mèo nhiều hơn số bạn thích Cá cảnh là
\choice
{$25$ bạn}
{\True $20$ bạn}
{$15$ bạn}
{$10$ bạn}
\loigiai{
$35 - 15 = 20$ bạn.
\textbf{Đáp án của câu hỏi là B}
}
\end{ex}

\textbf{Phần 2. Trắc nghiệm đúng sai}
\setcounterpageref{ex}{0}

\begin{ex} %[3X1V73]
Xét tính đúng sai của việc đọc và xử lý bảng số liệu:
\choiceTF
{\True Bảng số liệu giúp ta so sánh và tính toán các đại lượng dễ dàng hơn.}
{Các số liệu trong bảng luôn luôn phải được sắp xếp theo thứ tự tăng dần.}
{\True Từ bảng số liệu có thể tìm ra giá trị lớn nhất và nhỏ nhất.}
{Không thể tính tổng các số liệu trong cùng một cột của bảng.}
\loigiai{
\begin{itemize}
    \item Ý a) Đúng: Mục đích thống kê.
    \item Ý b) Sai: Sai vì có thể sắp xếp theo thời gian hoặc danh mục.
    \item Ý c) Đúng: Kĩ năng phân tích.
    \item Ý d) Sai: Sai vì ta có thể cộng tổng theo cột.
\end{itemize}
\begin{itemize}
    \item a) \textbf{Đúng}.
    \item b) \textbf{Sai}.
    \item c) \textbf{Đúng}.
    \item d) \textbf{Sai}.
\end{itemize}
}
\end{ex}

\begin{ex} %[3X1V73]
Số học sinh giỏi 4 khối lớp: Khối 1: 80 em; Khối 2: 95 em; Khối 3: 85 em; Khối 4: 90 em:
\choiceTF
{\True Khối 2 có nhiều học sinh giỏi nhất.}
{\True Khối 1 có ít học sinh giỏi nhất.}
{Tổng số học sinh giỏi của Khối 1 và Khối 3 bằng Khối 2.}
{\True Tổng số học sinh giỏi cả 4 khối là 350 em.}
\loigiai{
\begin{itemize}
    \item Ý a) Đúng: 95 em cao nhất.
    \item Ý b) Đúng: 80 em ít nhất.
    \item Ý c) Sai: Sai vì $80 + 85 = 165 \neq 95$.
    \item Ý d) Đúng: $80 + 95 + 85 + 90 = 350$ em.
\end{itemize}
\begin{itemize}
    \item a) \textbf{Đúng}.
    \item b) \textbf{Đúng}.
    \item c) \textbf{Sai}.
    \item d) \textbf{Đúng}.
\end{itemize}
}
\end{ex}

\textbf{Phần 3. Trả lời ngắn}
\setcounterpageref{ex}{0}

\begin{ex} %[3X1V73]
Số học sinh nam của 3 lớp lần lượt là: 18, 22, 20. Tính tổng số học sinh nam của cả 3 lớp.
\shortans{60}
\loigiai{
\textbf{Đáp số:} $60$\\
$18 + 22 + 20 = 60$ bạn.
}
\end{ex}

\begin{ex} %[3X1V73]
Tổ 1 gom được $45\text{ kg}$ giấy vụn, Tổ 2 gom được $55\text{ kg}$ giấy vụn. Cả hai tổ gom được bao nhiêu ki-lô-gam giấy vụn?
\shortans{100}
\loigiai{
\textbf{Đáp số:} $100$\\
$45 + 55 = 100\text{ (kg)}$.
}
\end{ex}

\begin{ex} %[3X1V73]
Bốn bạn quyên góp sách: An: 8 cuốn, Bình: 12 cuốn, Cúc: 10 cuốn, Dũng: 14 cuốn. Bạn góp nhiều nhất nhiều hơn bạn góp ít nhất bao nhiêu cuốn?
\shortans{6}
\loigiai{
\textbf{Đáp số:} $6$\\
$14 - 8 = 6$ cuốn.
}
\end{ex}

\begin{ex} %[3X1V73]
Số cây trồng được của lớp 3A trong 3 ngày: Ngày 1: 25 cây; Ngày 2: 30 cây; Ngày 3: 35 cây. Tính trung bình mỗi ngày lớp 3A trồng được bao nhiêu cây?
\shortans{30}
\loigiai{
\textbf{Đáp số:} $30$\\
$(25 + 30 + 35) : 3 = 90 : 3 = 30$ cây.
}
\end{ex}

\begin{ex} %[3X1V73]
Trong một hộp có 12 chiếc bút màu đỏ và 18 chiếc bút màu xanh. Hỏi số bút màu xanh nhiều hơn số bút màu đỏ bao nhiêu chiếc?
\shortans{6}
\loigiai{
\textbf{Đáp số:} $6$\\
$18 - 12 = 6$ chiếc.
}
\end{ex}

\begin{ex} %[3X1V73]
Một nông trại nuôi 3 chuồng lợn: chuồng A có 45 con, chuồng B có 55 con, chuồng C có 60 con. Hỏi chuồng C nhiều hơn chuồng A bao nhiêu con lợn?
\shortans{15}
\loigiai{
\textbf{Đáp số:} $15$\\
$60 - 45 = 15$ con.
}
\end{ex}
